\section{The Verification Framework with \Dargent}
{\label{sec:verification}}

%As we saw in~\S\ref{ss:formalised-semantics}, 
As we have demonstrated, \Dargent affects both
the formalisation of the type system and the C code generation of the compiler.
Moreover, it also affects the generated correspondence proof between
the C code and the pure shallow embedding in Isabelle/HOL. For this proof,
only the refinement between \Cogent's stateful update semantics
and the generated C code must change.
Above this low-level layer, the functional embeddings of \cogent still use
high-level algebraic types, regardless of the specified layouts.
%by the way data
% In addition to the extension of the type system, which impacts
% the deep embedding,
% the main component that must be adapted is the proof of
% correspondence between the update semantics and the C code. The crucial
% change is that getting or setting a field in \cogent does
% not correspond anymore to getting or setting a corresponding C field,
% but to calling a custom getter or setter.
%Indeed, at the C level, a field is now encoded in an array of words
%according to the layout.  Getting or setting a field in \cogent corresponds to
%a call to a generated custom getter or setter, in C.
In the following subsections,
we discuss  the extensions to the verification framework regarding:
%as \cogent records with layouts now compile to flat arrays:
%In the following subsections, we discuss what changes to the verification
%framework were necessary in this correspondence statement to take \Dargent into account:
\begin{enumerate}
  \item the correspondence between \cogent record values and C flat arrays;
  \item the correspondence between record operations in \cogent and C;
  \item the compositional properties of generated getters and setters;
  \item the correspondence between generated getters/setters and the specified
     layout.
\end{enumerate}
These extensions to the verification framework not only ensure that the C code
generated by the compiler is a refinement of the \cogent code, but also
prove that \cogent types are laid out in C correctly according to the programmer's \dargent
specification. Again, all the results we present in this section are
developed and checked in Isabelle/HOL.


The verification of any software system must assume the correctness of a 
\emph{trusted computing base}. In \cogent systems, this trusted computing base is largely
comprised of externally provided C code~\citep{OConnor:2016}. While the
\cogent framework allows for manual verification of this C 
code~\citep{Cheung_OR_22}, with \dargent we eliminate large swathes of this C code entirely, specifically the so-called ``glue code'' which translates between data formats. This reduces the size of the TCB without imposing any additional verification burden.
 


% In a last subsection, we describe additional generated proofs ensuring that C field getters and
% setters respect the given record layout specification.

% \review{what happens? what are the
%   implications? why not tackled? $\times$2.}

% \review{what proofs are currently generated, what proofs are still manual,
% what is still trusted not at all proven, and what will be done versus what will be left trusted. what guarantees Dargent currently provides and what could go wrong. which lemmas stated in the paper are formally proven. when something is not yet supported, what currently happens and the implications of that.}

% \review{are the lemmas formally proved? }

\subsection{Relating Record Values}\label{ssec:record-value-rel}

The correctness of compilation is justified by a \emph{refinement proof} that states that a simulation relation, called a \emph{refinement relation}, holds between the source language and the target language.
The refinement relation from the \cogent update semantics 
and the compiled C program is defined in terms of a relation between
\cogent values (in the sense of the update semantics)
and C values. 
Without \Dargent, this value
relation is straightforward for records: A \cogent record relates to a C structure with the same fields, such that each \cogent field is itself value-related to 
the corresponding C field.

With layouts, a \cogent record now relates to (essentially) a flat word array, and each
\cogent field relates to the result of applying the relevant custom
getter to the array.
This requires an appropriate embedding of our getters in Isabelle.
To this end, we have implemented an Isabelle/ML procedure that imports the C getters and setters generated for each field using the AutoCorres library, which produces a monadic, shallow HOL embedding that models the semantics of the C functions. Our ML procedure then simplifies these embeddings to pure Isabelle functions.
We call these simplified functions
\emph{direct getters} and \emph{direct setters}.
These direct getters and setters allow a simple statement
of the compositional properties between getters and setters,
detailed in~\autoref{sec:get-set-lemmas}.
%\subsubsection{Overview of the value relation}

% The correspondence statement between the deep embedding with update
% semantics and the AutoCorres embedding~\citep[Definition 2]{OConnor:2016} 
% relies on the definition of a value relation,
% for each C type, between:
% \begin{itemize}
%   \item a value, in the sense of the update semantics, and
%   \item an Isabelle term whose type represents the C type according to
%     AutoCorres.
% \end{itemize}
% When embedding C code in Isabelle using the AutoCorres library,
% each C type is represented directly as an Isabelle type. For example, the C structure
% \lstinline[language=C]${ U8 a; U16 b; }$ essentially embeds as
% an Isabelle record type consistings of two fields $a$ and $b$ with the types 
% \lstinline$8 word$ and \lstinline$16 word$ of natural numbers fitting on
% 8 and 16 bits respectively.
%  Then, C functions are embedded as \emph{monadic} Isabelle programs for the monad
%  \texttt{NonDet}, allowing for failure, non determinism, and stateful operations. 
% For example, a C
% function \texttt{U16\ square(U8\ a)} is embedded as a term of type
% \texttt{8 word $\Rightarrow$\ (16 word) NonDet} in Isabelle. A main motivation
% of \cogent is actually to get a pure non-monadic embedding.

% The value relation is defined automatically (by some ML procedure) for
% each involved C type in the generated C code. When no custom layout is
% involved, the value relation for a C struct between a value
% \texttt{v} and a term \texttt{t} of type \texttt{T} (representing the
% C type) enforces \texttt{v} to be a record with the same number of
% fields as \texttt{T}, such that the field values are themselves
% value-related to the corresponding fields of \texttt{t}. The adequacy of
% this generated value relation relies on the fact that records are
% straightforwardly compiled to C structures, when no layout is involved.

% \subsubsection{Records with custom layouts}

% When assigned a layout $\ell$, a \cogent record type
% %\texttt{T =  \{\ a\ :\ A,\ b\ :\ B\ \}}
% \lstinline$T = { a : A, b : B }$ translates (essentially) to an array of
% words \texttt{$\sem{ T }_\ell$}. The corresponding value relation between
% a value (in the sense of the update semantics) and a
%  word array \texttt{arr} still enforces the value to be a record type with the same number of fields as
% in \texttt{T}, and moreover that the field values relate
%  to the corresponding data encoded in the array
% of words according to the layout.

% In the particular example considered here, the value relation would enforce that the first
% field of the value relates to \texttt{get\_a\ arr}, where
% \texttt{get\_a} is the custom getter whose expected type is
% \texttt{$\sem{ T }_\ell$\ $\Rightarrow$ $\sem{ A }$}.

% \subsubsection{Direct definition of custom
% getters}
% \label{ss:direct-get}
% AutoCorres provides an Isabelle definition of the custom getters from
% the C code, but they do not have the right type for the value relation, as
% they are monadic. In the example above, it provides
% \texttt{get\_a\ :\ $\sem{ T }_\ell$\ $\Rightarrow$ $\sem{ A
%   }$\ NonDet}. Defining the value relation as described above thus requires 
% a \emph{direct} definition \texttt{direct\_get\_a\ :\ $\sem{ T
%   }_\ell$\ $\Rightarrow$ $\sem{ A }$} to be inferred. This is done by a ML procedure which
% inspects the monadic definition and simplifies it.

%There are other possible ways to achieve this goal. For example, the
%compiler could be extended to generate these direct getters in Isabelle.


\subsection{Relating Record Operations}\label{sec:record-lemmas}

Without \Dargent, records directly compile to C structures, so
it is trivial to relate
\cogent record operations to their compiled C versions
and show that they correspond.
With layouts, the correspondence proof becomes slightly more involved, since it relates
getting or setting fields in \cogent to calling custom getters or setters
in C. These correspondence statements are proved automatically by custom tactics we have developed 
which exploit the compositional properties of getters and setters, detailed in the next section,
as well as the relation between their monadic and direct embeddings.


\subsection{Verifying the Custom Getters and Setters}
\label{sec:get-set-lemmas}

To relate \cogent record operations and operations on a C array
$\texttt{arr}$ that models the record type, getters and setters are
expected to obey the following compositional laws schematically
summarised below. Given an $\texttt{arr}$ array that models a record
type:
\begin{description}[leftmargin=\parindent,labelindent=\parindent]
  \item[\text{[Roundtrip]}]
    If a \cogent value $x$ relates to a C value $v$,
    then it must also relate to the C value
    $\texttt{get\_a\ (set\_a\ arr\ $v$)}$.
    This is a weakening of the more intuitive roundtrip statement
\texttt{get\_a\ (set\_a\ arr $v$)\ =\ $v$}, which does not always hold:
As a counterexample, consider a boolean field. Since C does not natively provide such a 1-bit type,
booleans are typically represented using a single byte \lstinline{U8}. But a custom getter
for a boolean field would always return 0 or 1, picking 
the relevant bit as specified in the layout. Thus, we obtain a counter-example by taking $v=3$.
  \item[\text{[Frame]}]
    For distinct fields \texttt{a} and \texttt{b},
    \texttt{get\_a\ (set\_b\ arr\ v)} is equal to \texttt{get\_a\ arr}.
\end{description}
These statements are subject to typing constraints that we do not detail for brevity.
%Here, \texttt{direct\_get/set} refers to the direct definition of
%getters/setters mentioned in~\S\ref{ss:direct-get}, while \texttt{ac\_get/set} refers to
% the monadic AutoCorres embedding of them.
% Direct setters are inferred by a similar process as the direct getters.
% \lstinline${ sum : < A U8 | B U8 > }$ with a layout assigning the same location
% to the payloads of both alternatives. The custom setter takes as arguments a word array 
% and a structure with three fields: one for the tag value, one for the A-payload,
% and one for the B-payload (this is indeed how \cogent variant types are compiled). 
% It puts one payload or the other, depending on the tag value, at the 
% location specified by the layout. This is clearly not injective in the second argument, and thus
% there is no chance to get the intuitive strict equality.

% Another example where the equality does not hold is that of bit fields, involving integers with a non 
% standard width, for example $\ConsN{U29}$. 
% Such an integer type does not exist natively in C (neither in \cogent, see~\S\ref{ssec:bitfields}
% for an account of them using abstract types), but it can be represented
% as a $\ConsN{U32}$. However, getting and setting a field of this type with a 
% $\ConsN{U32}$ value results in its truncation on twenty-nine bits.

\subsection{Ensuring Getters and Setters Comply with Layouts}
\label{ss:sanity-checks}

As noted before, our compiler correctness certificate
does not require that our generated getters and setters actually respect the layout specification---only that they satisfy the compositional properties stated previously.
Thus, it remains to ensure that the generated getters and setters comply with data layouts specified by the user.

% Writing a generic specification of getters and setters
%once and for all, that does not need to be generated for each program (and carefully
%checked by the user), requires a deep embedding of them.
% AutoCorres relies under the hood on the C parser library that could be used for this purpose.
% However, we adopt another strategy:
To this end, we  generate specifications for getters in the deep embedding of \cogent,
returning \cogent values (in the update semantics) from a word array,
according to the layout.
More specifically, we implemented a
generic getter \lstinline|uval_from_array|
as a recursive function in Isabelle that takes as input a \cogent type, a
\Dargent layout,
a word array\footnote{Each of these inputs is of a suitable type inductively
  defined in Isabelle.},
and retrieves the value of the field of the given
\cogent type from the word array according to the given field layout.



We designed tactics to automatically show that generated getters respect this
specification. More precisely, given a field $b$ of \cogent type $\tau$ with associated layout
$\ell$, the tactics
prove that for any word array $a$ that has a valid
format according to the layout, \lstinline|get_b|~$a$
relates to \lstinline|uval_from_array|~$\tau\ \ell\ a$ in the sense of the value
relation,
 where \lstinline|get_b| is the 
direct getter.

We have not yet done the same for setters, although the compositional properties we verify,
together with the correctness of getters, already provide some formal guarantees. Additionally, we also automatically prove explicitly that
setters do not flip any bit that is outside the field layout. 

%\subsection{Tactics for the get/set lemmas}
%We wrote one ML tactic for each kind of get/set lemmas. They are about 50 lines long 
%each and rely on the tactics \texttt{word\_bitwise} and \texttt{monad\_eq} provided by the
%AutoCorres library, that simplify equalities between integers and monadic programs,
%respectively.
% Optimising them is left for future work.
% The first get/set lemma mentioned above
% can be unnecessarily slow to prove in the case where multiple variant types are involved,
% as the tactic currently checks every possible combination of constructors.

\subsection{Endianness}\label{ssec:endianness-verif}

Since \Cogent relies on the AutoCorres library, which does
not support big-endian architectures yet, we assume that the machine 
is little-endian.
Therefore, only big-endian annotations need to be seriously accounted for: 
Big-endian fields of type $\PrimType{U16}$, $\PrimType{U32}$, and $\PrimType{U64}$ require reversing the bytes on retrieval and setting. This reversal is performed in dedicated optimised functions.
As an example, here is the code for reversing the bytes in a $\PrimType{U32}$:
\begin{lstlisting}[language=C]
  static inline u32 swap_u32(u32 v) {
    v = ((v << 8) & 0xFF00FF00) | ((v >> 8) & 0xFF00FF);
    return (v << 16) | (v >> 16);
  }
\end{lstlisting}
This additional reversal step does not require any
extension to the proof tactics for the compositional properties of getters and
setters described in \autoref{sec:get-set-lemmas}, as they proceed first by unfolding all involved function definitions and applying an automated simplifier for bitwise operations, which can automatically prove that these swap operations are an involution. 
Recall that the value relation described in \autoref{ssec:record-value-rel} 
is defined based on our generated custom getters, which already take endianness into account, so the value relation also naturally takes 
endianness into account. 

 The proof described in \autoref{ss:sanity-checks}, that our getters respect the given layout, also must accommodate endianness. 
Specifically, we use a purely functional HOL embedding of the swap function to 
specify getters of big-endian values. Because our HOL embedding follows the
C implementation closely, our proof tactics also need no extension for big-endian fields to prove correctness of getters.

We also prove as a sanity check that these swap functions are correct, in the sense 
 that they reverse the byte order.

